\documentclass[a4paper,10pt]{article}

\usepackage[
    a4paper,
    top=12mm,
    bottom=12mm,
    left=12mm,
    right=12mm
]{geometry}

\usepackage{fontspec}
\usepackage{polyglossia}
\setdefaultlanguage[babelshorthands=true]{russian}
\setotherlanguage{english}

\setmainfont{FreeSerif}
\newfontfamily{\russianfonttt}[Scale=0.7]{DejaVuSansMono}

\usepackage[tiny, compact]{titlesec}
\usepackage{titling}
\setlength{\droptitle}{-1cm}

\usepackage{xurl}
\usepackage{hyperref}
\usepackage{bookmark}
\usepackage{array}
\usepackage{booktabs}
\usepackage{float}
\usepackage{csquotes}

\hypersetup{
    colorlinks=true,
    linkcolor=black,
    urlcolor=blue,
    citecolor=black,
    breaklinks=true
}

\begin{document}

\title{Анализ графов}
\author{\normalsize Пентегова Мария Сергеевна}
\date{}
\maketitle

\begin{flushright}
\begin{tabular}{r}
    Группа: \emph{\enquote{25.Б41-мм}} \\
    Кафедра: \emph{системного программирования} \\
    Научный руководитель: \emph{Григорьев Семён Вячеславович} \\
    Номер семестра практики: \emph{2} \\
\end{tabular}
\end{flushright}

\section{Контекст работы и постановка задачи}

Настоящая работа выполнена в рамках учебной практики на кафедре системного программирования. Исследование посвящено анализу эффективности алгоритмов обработки больших разреженных графов. Данная область критически важна для ускорения работы комплексных сетевых систем, социальных платформ и поисковых движков. 

Обход в ширину (\foreignlanguage{english}{Breadth-First Search, BFS})~--- это базовый алгоритм теории графов, при котором вершины посещаются последовательно по уровням удалённости от исходного источника. Скорость выполнения этого алгоритма определяет общую производительность систем интеллектуального анализа данных. Работа опирается на применение открытого программного интерфейса \texttt{GraphBLAS} (в реализации \texttt{SuiteSparse:GraphBLAS})~--- интерфейса прикладного программирования (\foreignlanguage{english}{Application Programming Interface, API}), который позволяет переформулировать графовые операции в терминах линейной алгебры над разреженными матрицами и векторами~\cite{graphblas-api}. 

В рамках исследования рассматриваются два вычислительных варианта алгоритма: \texttt{parent-BFS} (вариант с фиксацией дерева предков для каждой достижимой вершины) и \texttt{multisource BFS} (параллельный обход из фиксированного множества начальных источников).

\textbf{Целью работы является} экспериментальное исследование влияния способов представления структур графа и использования матрично-векторного интерфейса прикладного программирования \texttt{GraphBLAS} на время выполнения алгоритмов обхода в ширину.

\textbf{Актуальность и практическая польза} результатов заключаются в получении объективных критериев выбора архитектуры программного обеспечения для обработки больших данных. Эмпирические результаты позволяют определить границы применимости линейно-алгебраического подхода в зависимости от топологических свойств обрабатываемых сетей, что снижает непроизводительные затраты вычислительных ресурсов.

\textbf{Для достижения поставленной цели были сформулированы и решены следующие задачи}:
\begin{itemize}
    \item реализовать варианты \texttt{parent-BFS} и \texttt{multisource BFS} на языке C с использованием формата сжатого хранения строк (\foreignlanguage{english}{Compressed Sparse Row, CSR});
    \item разработать эквивалентные алгоритмы обхода средствами низкоуровневого интерфейса прикладного программирования библиотеки \texttt{SuiteSparse:GraphBLAS};
    \item спроектировать автоматизированный тестовый набор для проверки функциональной корректности разработанных программных модулей;
    \item измерить точное время выполнения алгоритмов на семи реальных грахах из репрезентативной коллекции разреженных матриц Университета Флориды (\foreignlanguage{english}{SuiteSparse Matrix Collection})~\cite{ssmc};
    \item выполнить детальный сравнительный анализ времени работы классического подходов на основе \texttt{CSR} и матрично-векторного подходов на основе \texttt{GraphBLAS}.
\end{itemize}

\section{Описание программного решения}

В рамках выполнения работы спроектирована и реализована программная система, разделенная на независимые вычислительные модули на языке C и подсистему верификации инвариантов.

В классических вариантах алгоритмов структура графа организуется в формате сжатого хранения строк \texttt{CSR}, минимизирующем объем занимаемой оперативной памяти. Смежные вершины располагаются непрерывно в одномерном массиве \texttt{col\_idx}. Для произвольной вершины $v$ её соседи находятся в строго детерминированном диапазоне памяти:
\[
[\texttt{row\_ptr}[v], \texttt{row\_ptr}[v+1])
\]
При первом обнаружении вершины $u$ из текущей вершины $v$ алгоритм \texttt{parent-BFS} выполняет запись в результирующий вектор: $\texttt{parent}[u] = v$. В алгоритме \texttt{multisource BFS} начальная очередь параллельно заполняется индексами всех стартовых вершин. Теоретическая асимптотическая сложность одного итерационного цикла обхода на структуре \texttt{CSR} составляет $O(|V|+|E|)$.

В версиях на основе \texttt{GraphBLAS} топология графа отображается в булеву матрицу смежности $A$, а текущий волновой фронт обхода~--- в разреженный булевый вектор $f_k$. Вычисление последующего фронта $f_{k+1}$ математически формулируется как операция матрично-векторного умножения:
\[
f_{k+1} = f_k A
\]
Исключение повторных посещений пройденных вершин аппаратно-независимо оптимизируется за счет применения встроенного в интерфейс прикладного программирования механизма маскирования. 

Для верификации разработанного программного обеспечения был спроектирован автоматизированный тестовый набор, обеспечивающий полное покрытие кодовой базы. Тестирование завершилось со стопроцентным успешным результатом (\texttt{100\% passed}). Инфраструктура проверки включает два независимых этапа:
\begin{itemize}
    \item \textbf{Тестирование инвариантов.} На фиксированных синтетических топологиях автоматически проверяются метрики достижимости, глубина вычисленных уровней, связность дерева предков и строгое отсутствие циклических зависимостей;
    \item \textbf{Перекрёстная валидация.} Результаты явного последовательного обхода структуры \texttt{CSR} поэлементно сравниваются с векторами, полученными в результате вычислений \texttt{GraphBLAS} при эквивалентных стартовых условиях.
\end{itemize}

Для обеспечения строго контролируемых условий эксперимента все исследуемые алгоритмы обхода были разработаны самостоятельно непосредственно поверх базового интерфейса на языке C библиотеки \texttt{SuiteSparse:GraphBLAS}. Использование готовых высокоуровневых реализаций из сторонних проектов, таких как \texttt{LAGraph} или библиотека обобщённых алгоритмов на графах \texttt{Boost Graph Library}~\cite{lagraph,boost-graph}, было намеренно исключено для обеспечения изоляции измеряемых факторов. Готовые библиотеки содержат внутренние недокументированные эвристики диспетчеризации, автоматическое динамическое переключение между разреженными и плотными векторами фронта, а также встроенные высокоуровневые оптимизации управления памятью. Включение таких абстракций в контур тестирования не позволило бы провести изоляцию факторов и корректно сопоставить базовую эффективность самой матрично-векторной концепции \texttt{GraphBLAS} с разработанным низкоуровневым последовательным обходом структуры \texttt{CSR}. Самостоятельная реализация гарантирует применение единого распределителя памяти, идентичные критерии останова по пустому фронту и эквивалентные стартовые условия запуска для всех четырёх рассматриваемых вариантов алгоритма.


\section{Методика и анализ результатов экспериментального исследования}

Экспериментальное исследование направлено на измерение точного времени выполнения алгоритмов. В качестве измеряемой метрики принято среднее арифметическое времени работы в миллисекундах. Все эксперименты проводились на изолированном вычислительном стенде со следующими фиксированными параметрами окружения:
\begin{itemize}
    \item центральный процессор Intel Core Ultra 7 255H;
    \item объем доступной оперативной памяти составляет 30.000 ГБ;
    \item операционная система Ubuntu 26.04.1 с длительным периодом поддержки (\foreignlanguage{english}{Long Term Support, LTS});
    \item набор компиляторов \texttt{GCC} (\foreignlanguage{english}{GNU Compiler Collection}) версии 15.2.0;
    \item библиотека SuiteSparse:GraphBLAS версии 7.4.0;
    \item параметры оптимизации компилятора: \texttt{-O3 -DNDEBUG} (стандарт C11).
\end{itemize}

В качестве входных данных использовались семь разреженных матриц из коллекции \texttt{SuiteSparse Matrix Collection}~\cite{ssmc}. Загрузка осуществлялась из файлов матричного формата Matrix Market. Матрицы с флагом \texttt{general} интерпретировались как ориентированные графы, а с флагом \texttt{symmetric}~--- как неориентированные. При конвертации выполнялось только построение \texttt{CSR}-структур без изменения направлений дуг. Выбор стартовой вершины для одиночного поиска производился по критерию максимальной исходящей степени (\texttt{out-degree}), а для мультипоиска фиксировались четыре индекса: $0$, $n/4$, $n/2$ и $3n/4$. Характеристики исследованных графов приведены в таблице~\ref{tab:graphs}.

\begin{table}[H]
\centering
\scriptsize
\setlength{\tabcolsep}{6pt}
\begin{tabular}{lrr}
\toprule
Граф & Число вершин $|V|$, млн & Число рёбер $|E|$, млн \\
\midrule
\texttt{cit-Patents} & 23.947 & 28.854 \\
\texttt{com-Amazon} & 0.335 & 0.926 \\
\texttt{com-Orkut} & 3.072 & 117.185 \\
\texttt{roadNet-CA} & 0.916 & 5.105 \\
\texttt{road\_usa} & 3.775 & 16.519 \\
\texttt{soc-LiveJournal1} & 4.848 & 68.994 \\
\texttt{web-Google} & 1.971 & 2.767 \\
\bottomrule
\end{tabular}
\caption{Характеристики графов, использованных в книге экспериментов.}
\label{tab:graphs}
\end{table}

Для минимизации случайных погрешностей операционной системы для каждой программной конфигурации выполнялось по 10 независимых изолированных запусков. Время предварительного выделения памяти и построения внутренних объектов \texttt{GraphBLAS} было полностью исключено из замеров, что позволило оценивать исключительно чистое время выполнения обхода. Все численные значения в таблице~\ref{tab:results} приведены к единой точности~--- три знака после запятой.

\begin{table}[H]
\centering
\scriptsize
\setlength{\tabcolsep}{5pt}
\begin{tabular}{lrrrr}
\toprule
Граф & \texttt{CP}, мс & \texttt{CM}, мс & \texttt{GB-L}, мс & \texttt{GB-M}, мс \\
\midrule
\texttt{cit-Patents} & 1034.142 & 1160.571 & 26347.954 & 24488.670 \\
\texttt{com-Amazon} & 29.881 & 22.404 & 62.855 & 61.161 \\
\texttt{com-Orkut} & 2121.789 & 1993.293 & 542.795 & 582.770 \\
\texttt{roadNet-CA} & 52.984 & 48.788 & 97.838 & 109.488 \\
\texttt{road\_usa} & 7.797 & 0.491 & 189.076 & 36.194 \\
\texttt{soc-LiveJournal1} & 1613.440 & 1662.308 & 439.515 & 390.446 \\
\texttt{web-Google} & 88.148 & 79.361 & 713.218 & 688.762 \\
\bottomrule
\end{tabular}
\caption{Среднее время выполнения BFS по 10 запускам (мс), где \texttt{CP}~--- классический \texttt{parent-BFS}, \texttt{CM}~--- классический \texttt{multisource BFS}, \texttt{GB-L}~--- GraphBLAS BFS с вычислением уровней, \texttt{GB-M}~--- GraphBLAS Multisource BFS.}
\label{tab:results}
\end{table}

Анализ экспериментальных данных указывает на сильную зависимость эффективности подходов от топологической структуры графа:
\begin{itemize}
    \item \texttt{GraphBLAS}-реализации (\texttt{GB-L} и \texttt{GB-M}) продемонстрировали существенное преимущество на плотных социальных грахах с выраженной кластеризацией (\texttt{com-Orkut} и \texttt{soc-LiveJournal1}). На графе \texttt{com-Orkut} среднее время выполнения сократилось примерно в 3.908 раза, на \texttt{soc-LiveJournal1}~--- в 3.670 раза. Это объясняется высокой эффективностью внутренних векторных оптимизаций разреженных матричных вычислений при обработке сверхкрупных фронтов с высокой степенью связности вершин;
    \item Классический \texttt{CSR}-подход (\texttt{CP} и \texttt{CM}) оказался значительно эффективнее на разреженных сетях с большим топологическим диаметром и низкой средней степенью вершин~--- дорожных грахах (\texttt{roadNet-CA}, \texttt{road\_usa}) и иерархических структурах (\texttt{cit-Patents}, \texttt{web-Google}). На графе \texttt{road\_usa} для мультипоиска последовательный обход превзошёл матрично-векторный аналог в 73.714 раза. Накладные расходы на диспетчеризацию вычислений внутри ядра \texttt{GraphBLAS} на таких структурах превышают время полезной работы.
\end{itemize}

Зафиксированные показатели принимаются как верифицированный эмпирический базис исследования. Поуровневый замер динамического изменения размера фронтов в рамках текущей работы не производился.

\section{Заключение}

В ходе выполнения работы получены следующие результаты:
\begin{itemize}
    \item на языке C успешно реализованы 4 варианты алгоритма обхода в ширину с использованием различных подходов к организации структуры данных (\texttt{CSR} и матрично-векторный интерфейс прикладного программирования \texttt{GraphBLAS});
    \item разработан автоматизированный тестовый набор, подтвердивший логическую корректность кода и выполнение всех проверяемых инвариантов (\texttt{100\% passed});
    \item проведены комплексные замеры времени работы алгоритмов на 7 графах из коллекции \texttt{SuiteSparse}, позволившие выявить границы применимости исследуемых технологий;
    \item эмпирически доказана высокая эффективность применения \texttt{GraphBLAS} на социальных сетях и безусловное преимущество низкоуровневых структур хранения \texttt{CSR} на дорожных графах.
\end{itemize}

Репозиторий проекта: \url{https://github.com/MariaPentegova/graph_blas_bfs-2}. 

Исходный код вычислений, тесты, конфигурации автоматической сборки проекта, а также файлы матриц, размещённые с применением расширения системы контроля версий для работы со сверхкрупными файлами \texttt{Git Large File Storage (Git LFS)}~\cite{git-lfs}, доступны в основном каталоге репозитория.

\begin{thebibliography}{8}

\setlength{\itemsep}{0.05em}
\setlength{\parskip}{0pt}

\bibitem{graphblas-api}
GraphBLAS Forum.
\textit{The GraphBLAS C API Specification. Version 2.1}, 2023.
\hfill \\URL: \url{https://graphblas.org/docs/GraphBLAS_API_C_v2.1.0.pdf}
(дата обращения: 24.09.2026).

\bibitem{graphblas}
Davis, T. A.
\textit{SuiteSparse:GraphBLAS}.
Официальный репозиторий проекта.
\hfill \\URL: \url{https://github.com/DrTimothyAldenDavis/GraphBLAS}
(дата обращения: 24.09.2026).

\bibitem{lagraph}
GraphBLAS.org.
\textit{LAGraph}.
Официальный репозиторий проекта.
\hfill \\URL: \url{https://github.com/GraphBLAS/LAGraph}
(дата обращения: 24.09.2026).

\bibitem{boost-graph}
Boost.
\textit{Boost Graph Library}.
Официальная документация.
\hfill \\URL: \url{https://www.boost.org/doc/libs/latest/libs/graph/doc/html/graph/index.html}
(дата обращения: 24.09.2026).

\bibitem{ssmc}
Davis, T. A., Hu, Y.
\textit{The University of Florida Sparse Matrix Collection}.
ACM Transactions on Mathematical Software, 38(1), 2011.
\hfill \\URL: \url{https://sparse.tamu.edu/}
(дата обращения: 24.09.2026).

\bibitem{cmake-ctest}
Kitware.
\textit{CTest Documentation}.
\hfill \\URL: \url{https://cmake.org/cmake/help/latest/manual/ctest.1.html}
(дата обращения: 24.09.2026).

\bibitem{git-lfs}
Git LFS.
\textit{Git Large File Storage}.
\hfill \\URL: \url{https://git-lfs.com/}
(дата обращения: 24.09.2026).

\end{thebibliography}

\end{document}
